\chapter{Promedios y sumación de Cesàro}\label{cap:v4-01}
Una serie de Fourier puede conservar información sobre una función aunque sus sumas parciales no converjan. Antes de trabajar con funciones, estudiaremos una operación elemental que reduce las oscilaciones: tomar promedios. Conviene distinguir siempre los términos de una serie, sus sumas parciales y los promedios de esas sumas.

\section{Regularidad del promedio}\label{lectura:v4-01}
\index{Cesaro@Cesàro!sumacion@sumación}\index{serie!sumable por Cesaro@sumable por Cesàro}
\begin{proposicion}
Si $x_n\to x$ en un espacio normado $V$, entonces $n^{-1}\sum_{k=1}^nx_k\to x$.
\end{proposicion}
\begin{proof}
Fijemos $\eps>0$. Existe $N$ tal que $\norm{x_k-x}<\eps/2$ para $k>N$. Para $n>N$,
\[
 \left\|\frac1n\sum_{k=1}^nx_k-x\right\|
 \leq\frac1n\sum_{k=1}^N\norm{x_k-x}+\frac{n-N}{n}\frac\eps2.
\]
La primera suma tiene un número fijo de términos; su cociente entre $n$ es menor que $\eps/2$ cuando $n$ es suficientemente grande. No se ha usado completitud: el límite ya pertenece a $V$.
\end{proof}
\begin{definicion}
Para una serie $\sum_{n\geq1}v_n$ escribimos
\[
 s_n=\sum_{k=1}^nv_k,\qquad \sigma_n=\frac1n\sum_{k=1}^ns_k.
\]
La serie es \emph{sumable por Cesàro}, o $(C,1)$-sumable, con suma $v$, si $\sigma_n\to v$ en norma.
\end{definicion}
\begin{proposicion}
La sumación de Cesàro es lineal y consistente con la suma ordinaria. Además, si $\sigma_n$ converge, entonces $s_n/n\to0$.
\end{proposicion}
\begin{proof}
Para las series $\sum u_n$ y $\sum v_n$, las sumas parciales y sus promedios se combinan con los mismos escalares; la continuidad de suma y multiplicación escalar prueba linealidad. La proposición anterior, aplicada a $(s_n)$, prueba consistencia. Finalmente,
\[
 \frac{s_n}{n}=\sigma_n-\frac{n-1}{n}\sigma_{n-1}
 =(\sigma_n-\sigma_{n-1})+\frac{\sigma_{n-1}}n\quad(n\geq2).
\]
La primera diferencia tiende a cero y la segunda también, porque una sucesión convergente es acotada.
\end{proof}
\begin{ejemplo}
Si $z\in\C$, $|z|=1$ y $z\ne1$, entonces $s_n=z(1-z^n)/(1-z)$. En consecuencia,
\[
 \sigma_n=\frac z{1-z}-\frac{z^2(1-z^n)}{n(1-z)^2}\longrightarrow\frac z{1-z}.
\]
La serie $\sum z^n$ no converge ordinariamente, pues sus términos no tienden a cero. Esto incluye $z=e^{i}$, ejemplo de la lectura original. Los promedios no cambian la suma de una serie convergente, pero sí amplían la clase de series a las que puede asignarse una suma.
\end{ejemplo}

\section{Una condición tauberiana}
\index{Hardy--Landau, teorema de}\index{teorema!tauberiano}
\begin{teorema}[Hardy--Landau]
Si $\sigma_n\to v$ y $\norm{v_n}\leq A/n$ para todo $n$ suficientemente grande, entonces $s_n\to v$.
\end{teorema}
\begin{estrategia}
El promedio global $\sigma_n$ puede ocultar oscilaciones. La cota sobre $v_n$ asegura que las sumas parciales cambian poco dentro de un intervalo corto de índices. Compararemos $s_m$ con el promedio de $s_{n+1},\ldots,s_m$, tomando $m/n$ cercano a uno.
\end{estrategia}
\begin{proof}
Para $m>n$ después del índice donde vale la hipótesis,
\[
 \frac1{m-n}\sum_{k=n+1}^ms_k=\frac{m\sigma_m-n\sigma_n}{m-n}.
\]
Además, $\norm{s_m-s_k}\leq\sum_{j=k+1}^m A/j\leq A(m-k)/n$. Por tanto
\[
 \norm{s_m-v}\leq
 \frac{m\norm{\sigma_m-v}+n\norm{\sigma_n-v}}{m-n}
 +A\frac{m-n}{n}.
\]
Fijemos $0<\delta<1$ y, para cada $m$ grande, tomemos $n=\lfloor m/(1+\delta)\rfloor$. Entonces $n\to\infty$, $m/n\to1+\delta$ y los cocientes $m/(m-n)$ y $n/(m-n)$ permanecen acotados. El primer sumando tiende a cero, de modo que $\limsup_m\norm{s_m-v}\leq A\delta$. Como esto vale para todo $\delta>0$, el límite superior es cero.
\end{proof}
Esta prueba corrige las identidades de sumas intermedias de la lectura: no aparece ningún término adicional $s_{n+1}$ en $\sum_{k=n+1}^ms_k=m\sigma_m-n\sigma_n$.

\section{Ejercicios sobre sumabilidad y productos}
\begin{ejercicio}[Lectura 1, ejercicio 1]\label{ej:v4-l01-1}
Si $\sum v_n$ tiene suma de Cesàro $v$, demostrar que $\sum_{n\geq1}v_{n+1}$ tiene suma de Cesàro $v-v_1$.
\end{ejercicio}


\begin{ejercicio}[Lectura 1, ejercicio 2: rectificación del criterio de Abel]\label{ej:v4-l01-2}
El enunciado original afirma que la monotonía de $(\alpha_n)$ basta para que $\sum\alpha_nv_n$ sea sumable por Cesàro cuando lo es $\sum v_n$. Mostrar que es falso. Probar la siguiente versión suficiente en un espacio de Banach $V$: si $(\alpha_n)$ converge y
\[
 \sum_{n\geq1} n|\alpha_n-\alpha_{n+1}|<\infty,
\]
entonces el producto término a término conserva la sumabilidad de Cesàro.
\end{ejercicio}


\begin{ejercicio}[Lectura 1, ejercicio 3]\label{ej:v4-l01-3}
Sean $s_n=\sum_{k=1}^nv_k$, $t_n=\sum_{k=1}^nkv_k$ y $\sigma_n=n^{-1}\sum_{k=1}^ns_k$. Probar: (i) $t_n=(n+1)s_n-n\sigma_n$; (ii) si la serie es Cesàro-sumable, converge ordinariamente si y solo si $t_n/n\to0$; (iii) la serie es Cesàro-sumable si y solo si $\sum t_n/[n(n+1)]$ converge en $V$.
\end{ejercicio}


\begin{ejercicio}[Lectura 1, ejercicio 4]\label{ej:v4-l01-4}
En dimensión finita, demostrar que una serie es Cesàro-sumable si y solo si lo son sus series coordenadas respecto de una base fija.
\end{ejercicio}


\begin{ejercicio}[Lectura 1, ejercicio 5: uniformidad necesaria]\label{ej:v4-l01-5}
El original propone caracterizar la sumabilidad en $\ell^\infty$ por sumabilidad de cada coordenada y acotación de las sumas obtenidas. Dar un contraejemplo y formular el criterio correcto.
\end{ejercicio}


\begin{ejercicio}[Lectura 1, ejercicio 6]\label{ej:v4-l01-6}
En un álgebra de Banach $A$, sea $\sum u_n$ absolutamente convergente y sea $w_n=\sum_{k=1}^nu_kv_{n+1-k}$. Demostrar que $\sum w_n$ converge absolutamente, ordinariamente o por Cesàro si $\sum v_n$ tiene, respectivamente, esa propiedad. En cada caso su suma es el producto de las sumas en el orden indicado.
\end{ejercicio}


\begin{ejercicio}[Lectura 1, ejercicio 7]\label{ej:v4-l01-7}
Si $\sum u_n=a$ y $\sum v_n=b$ convergen en un álgebra normada, probar que su producto de Cauchy tiene suma de Cesàro $ab$. Determinar si necesariamente converge ordinariamente.
\end{ejercicio}


\begin{ejercicio}[Lectura 1, ejercicio 8]\label{ej:v4-l01-8}
Calcular la suma de Cesàro de $\sum_{n\geq1}(-1)^{n+1}$ y mostrar que su producto de Cauchy consigo misma no es Cesàro-sumable.
\end{ejercicio}


\paragraph{Enlace con las series de Fourier.}
Los promedios de Cesàro convertirán el núcleo oscilante de Dirichlet en el núcleo no negativo de Fejér. La positividad permitirá estimar el error separando una vecindad pequeña del origen de su complemento; esa idea reaparecerá en las identidades aproximadas de grupos.
